\chapter{Background and Literature Review}\label{ch:background-and-literature-review}

\section{Dynamically Balanced Mobile Robots}\label{dynamically-balanced-mobile-robots}

\subsection{Wheeled inverted-pendulum dynamics}\label{wheeled-inverted-pendulum-dynamics}

A self-balancing wheeled robot is commonly modeled as an inverted
pendulum whose pivot is accelerated by the driven wheels. The upright
configuration is open-loop unstable: a small pitch displacement produces
a gravitational moment that increases the displacement unless wheel
acceleration moves the contact region beneath the center of mass.
Grasser et al. demonstrated this coupling in JOE, a two-wheel mobile
inverted pendulum whose mechanical, power-electronic, and control
subsystems were designed together \cite{grasser2002joe}. Pathak, Franch, and Agrawal
subsequently used partial feedback linearization to separate velocity
and position objectives from the internal balance dynamics \cite{pathak2005wip}.
These systems establish the central lesson for Zephyr: translation and
pitch cannot be designed as independent subsystems when the same wheel
torques support both.

The full equations are nonlinear because the body center of mass moves
on an arc about the wheel axis and because motor reaction torque acts on
the body. Linear control is nevertheless useful near the upright
equilibrium. A linearized state-space model permits controllability
analysis, pole placement, and linear-quadratic regulation. LQR selects a
state-feedback gain by minimizing a quadratic measure of state deviation
and control effort \cite{anderson1990optimal}, while nonlinear analysis is still required
to understand the finite region in which the linear approximation is
credible \cite{khalil2002nonlinear}. Recent comparative work continues to use LQR as a
transparent baseline for inverted-pendulum stabilization \cite{chacko2023trajectory}. For
Zephyr, a linear controller is therefore appropriate as an initial
implementation if the nonlinear plant, actuator saturation, and fall
thresholds remain part of the validation.

\subsection{Underactuation, nonminimum-phase behavior, and practical limits}\label{underactuation-nonminimum-phase-behavior-and-practical-limits}

A wheel torque both accelerates the base and applies an equal and
opposite reaction torque to the body. The balance coordinate is not
independently actuated, which makes the robot underactuated. A commanded
translation generally requires a transient body lean, and an aggressive
position loop can conflict with the inner balance loop. This behavior
motivates a cascade in which pitch stabilization has the highest
bandwidth, forward velocity or position is slower, and yaw and lateral
objectives are allocated without consuming the torque reserve required
for balance. The final bandwidth separation must be based on measured
delay and actuator dynamics rather than a generic ratio.

Practical limits define the recoverable set. Finite motor current,
wheel-ground friction, gearbox backlash, battery voltage, sensor range,
and computation delay limit the pitch angle and angular rate from which
the robot can recover. A controller that is stable in an unconstrained
model may fail after a command saturates. Accordingly, the Zephyr model
includes torque and rate limits, and the proposed tests record both
pitch response and motor effort. No recovery envelope is claimed until
the physical robot is tested under controlled conditions.

\section{Omnidirectional and Mecanum-Wheel Mobility}\label{omnidirectional-and-mecanum-wheel-mobility}

\subsection{Mecanum-wheel principle}\label{mecanum-wheel-principle}

Ilon\textquotesingle s wheel patent describes a load-bearing wheel with
peripheral rollers arranged obliquely to the wheel plane \cite{ilon1975mecanum}.
Rotation of the main wheel creates a contact velocity that can be
decomposed into a constrained component and a roller-supported
component. Combining wheels of opposite handedness enables a platform to
synthesize motion directions that a conventional fixed wheel cannot. Pin
and Killough placed this concept within the broader family of holonomic
and omnidirectional platforms \cite{pin1994omnidirectional}, while Ferrière, Raucent, and
Campion examined the geometry and design of omnidirectional wheels
\cite{ferriere1996omnimobile}. Gfrerrer later provided a geometric treatment of the Mecanum
wheel itself \cite{gfrerrer2008geometry}.

The ideal kinematic model treats each roller as imposing a linear
velocity constraint. This abstraction is valuable for deriving a
Jacobian and checking controllability, but it does not fully represent
discrete roller contact, varying normal load, deformation, or slip.
Bayar and Ozturk showed that roller contact forces affect Mecanum-wheel
behavior \cite{bayar2020contact}. More recent experimental modeling work likewise
reports the need to consider contact and slip when high accuracy is
required \cite{bhargava2026review}, \cite{barcenas2025robust}. Zephyr is especially sensitive to these
effects because the wheel forces must stabilize pitch as well as produce
planar motion.

\subsection{Mobility classification and geometry rank}\label{mobility-classification-and-geometry-rank}

Campion, Bastin, and D\textquotesingle Andréa-Novel classified
wheeled-mobile-robot models according to their kinematic constraints and
mobility \cite{campion1996structural}. Muir and Neuman also showed how wheel constraints can
be assembled into a systematic kinematic model \cite{muir1987kinematic}. These
approaches emphasize that mobility follows from the constraint matrix,
not from a wheel label. A four-wheel Mecanum platform is not
automatically holonomic if the roller signs or locations make the rows
of its Jacobian dependent.

Indiveri derived kinematic and control relationships for Swedish-wheel
omnidirectional robots \cite{indiveri2009swedish}. Conventional four-corner formulas use
two geometric lever arms and a prescribed X or O roller arrangement.
They cannot be copied directly to Zephyr because all wheel centers lie
on one line. For a collinear arrangement, yaw arises from longitudinal
force differences across the axial wheel locations, while lateral force
arises from the roller signs. Full planar mobility requires both roller
handednesses and a distribution of wheel locations that keeps the
longitudinal, lateral, and yaw columns independent. Chapter 6 expresses
this requirement as a rank condition that can be evaluated after
measurement.

\begin{longtable}[]{@{}
  >{\raggedright\arraybackslash}p{(\columnwidth - 0\tabcolsep) * \real{1.0000}}@{}}
\toprule\noalign{}
\begin{minipage}[b]{\linewidth}\raggedright
\textless INSERT FIGURE 2.1: Create a synthesis diagram showing how
inverted-pendulum dynamics, Mecanum-wheel constraints, IMU state
estimation, actuator\slash{}CAN limits, real-time scheduling, and safety
supervision interact in Zephyr\textgreater{}\\
What the figure should show: Create a synthesis diagram showing how
inverted-pendulum dynamics, Mecanum-wheel constraints, IMU state
estimation, actuator\slash{}CAN limits, real-time scheduling, and safety
supervision interact in Zephyr\\
Likely source: author-created vector diagram based on references \cite{grasser2002joe}
through \cite{liu1973scheduling} and the repository architecture\strut
\end{minipage} \\
\midrule\noalign{}
\endhead
\bottomrule\noalign{}
\endlastfoot
\end{longtable}

\begin{center}\singlespacing
\phantomsection\addcontentsline{lof}{figure}{\protect\numberline{2.1}Relationship among balancing, omnidirectional mobility, state estimation, and embedded implementation.}
\textbf{Figure 2.1. Relationship among balancing, omnidirectional mobility, state estimation, and embedded implementation.}
\end{center}

\section{Control and State Estimation}\label{control-and-state-estimation}

\subsection{PID, state feedback, and LQR}\label{pid-state-feedback-and-lqr}

PID control remains attractive because it is easy to implement and tune
from measured responses. A pitch-angle loop can command wheel torque
from angle, rate, and integral error, while an outer velocity loop
shifts the pitch reference. The method becomes harder to reason about
when translation, pitch, yaw, lateral motion, actuator saturation, and
allocation interact. Anti-windup is essential whenever the command is
limited.

State feedback provides a direct way to use pitch, pitch rate,
translational velocity, and position in one control law. LQR supplies a
repeatable gain-selection framework once the linear model and weighting
matrices are documented \cite{astrom2008feedback}, \cite{anderson1990optimal}. It does not remove the need
for engineering judgment: state scaling changes the meaning of the
weights, unmodeled delay can erode phase margin, and actuator limits
invalidate the unconstrained optimum. The Zephyr simulation draft uses
LQR as a proposed design method, but no gain is described as implemented
until it appears in firmware and is verified on hardware.

\subsection{Inertial sensing and complementary filtering}\label{inertial-sensing-and-complementary-filtering}

An accelerometer provides a gravity-referenced attitude estimate when
translational acceleration is small, whereas a gyroscope measures
angular rate with high short-term fidelity but accumulates bias when
integrated. A complementary filter combines the low-frequency
accelerometer reference with the high-frequency integrated gyroscope
signal. Mahony, Hamel, and Pflimlin placed complementary attitude
estimation on a rigorous nonlinear foundation \cite{mahony2008filters}. The current
Zephyr code path includes an IMU driver, mounting transform, gyro
calibration, an attitude estimator, and an IMU state snapshot. The
physical mounting transform and sign convention must be verified by
controlled rotations before any controller is armed.

\section{Embedded Actuation and Communication}\label{embedded-actuation-and-communication}

\subsection{Integrated actuators and local motor control}\label{integrated-actuators-and-local-motor-control}

Zephyr uses four CubeMars AK40-10 actuators. The manufacturer describes
the unit as a brushless motor, 10:1 planetary gearbox, internal driver,
and magnetic encoder integrated into one module \cite{cubemars2026ak4010}. Published
ratings include 24 V nominal voltage, 1.3 N\ensuremath{\cdot}m rated torque, 4.1 N\ensuremath{\cdot}m peak
torque, 370 r\slash{}min rated speed, 435 r\slash{}min no-load speed, 2.7 A rated
current, and 7.3 A peak current. These are component ratings, not
selected operating limits. The installed firmware version, operating
mode, current ceiling, speed limits, and protection settings require
hardware confirmation.

An integrated actuator moves the inner current and
field-oriented-control loops out of the central MCU, reducing the
required update rate of the supervisory controller. It also hides some
motor dynamics and introduces protocol, firmware, and mode-dependent
behavior. For initial simulation, the actuator can be represented as a
saturated command path with a measured delay or first-order response. A
higher-order electrical motor model is justified only if current-loop
data and parameters are available. Gearbox backlash is relevant because
balance commands reverse near zero speed, where the wheel may not
immediately follow the internal motor encoder.

\subsection{CAN networks and real-time scheduling}\label{can-networks-and-real-time-scheduling}

CAN provides prioritized, frame-based communication and automatic error
detection. The ESP32 exposes a Classical CAN-compatible controller named
TWAI, but it still requires an external physical-layer transceiver
\cite{espressif2026esp32}, \cite{espressif2026twai}. The current prototype uses an Adafruit CAN Pal based
on the TJA1051T\slash{}3 transceiver \cite{nxp2026tja1051}, \cite{adafruit2026canpal}. A linear bus with
short stubs and one 120 \ensuremath{\Omega} termination at each physical end is the
baseline topology; TI\textquotesingle s selectable-termination reference
design explains why the distal terminations are required \cite{ti2026tida01238}. Bus
bitrate alone does not establish real-time performance. Frame
identifiers, payloads, feedback periods, arbitration, timeouts, and
error counters must be logged under load.

Real-time software must complete sensing, estimation, control, and
command transmission before their deadlines. Liu and Layland established
foundational schedulability results for periodic fixed-priority tasks
\cite{liu1973scheduling}. Zephyr uses FreeRTOS concepts, but its final schedule should
be verified by measurement because driver blocking, serial bursts,
display updates, logging, and CAN retries can create execution-time
variability \cite{freertos2026scheduling}. The safety architecture should treat stale data
and missed deadlines as faults rather than allowing the last command to
persist indefinitely.

\section{Research Gap and Design Implications}\label{research-gap-and-design-implications}

The literature separately provides mature methods for wheeled inverted
pendulums and for statically stable omnidirectional platforms. The
Zephyr concept places those methods in direct interaction. Its four
wheels add actuation redundancy, but the collinear geometry removes the
conventional rectangular footprint and changes the yaw moment arm. Its
Mecanum rollers provide lateral authority, but roller slip perturbs both
odometry and the force available for stabilization. Its integrated motor
controllers reduce central computation, but CAN delay and gearbox
compliance become part of the balance loop.

The practical gap addressed by this thesis is therefore a traceable
system design and validation method for a collinear four-wheel Mecanum
balancing robot. The work does not claim that no related arrangement
exists. Instead, it asks whether the actual geometry has full planar
rank, whether the built platform can stabilize under its measured
limits, and how closely a reduced model predicts that behavior. This
framing leads directly to the requirements and evidence matrix in
Chapter 3.

\par\medskip\noindent\singlespacing
\phantomsection\addcontentsline{lot}{table}{\protect\numberline{2.1}Comparison of relevant mobile-robot architectures and methods}
\textbf{Table 2.1. Comparison of relevant mobile-robot architectures and methods}\par

\begin{longtable}[]{@{}
  >{\raggedright\arraybackslash}p{(\columnwidth - 6\tabcolsep) * \real{0.2250}}
  >{\raggedright\arraybackslash}p{(\columnwidth - 6\tabcolsep) * \real{0.2667}}
  >{\raggedright\arraybackslash}p{(\columnwidth - 6\tabcolsep) * \real{0.2583}}
  >{\raggedright\arraybackslash}p{(\columnwidth - 6\tabcolsep) * \real{0.2500}}@{}}
\toprule\noalign{}
\begin{minipage}[b]{\linewidth}\raggedright
\textbf{System or method}
\end{minipage} & \begin{minipage}[b]{\linewidth}\raggedright
\textbf{Established contribution}
\end{minipage} & \begin{minipage}[b]{\linewidth}\raggedright
\textbf{Limitation relative to Zephyr}
\end{minipage} & \begin{minipage}[b]{\linewidth}\raggedright
\textbf{Design implication}
\end{minipage} \\
\midrule\noalign{}
\endhead
\bottomrule\noalign{}
\endlastfoot
Two-wheel inverted pendulum \cite{grasser2002joe}, \cite{pathak2005wip} & Coupled translation and
pitch dynamics; balance and position control & No direct lateral-motion
allocation & Preserve a high-bandwidth balance objective while adding
redundant wheel allocation. \\
Four-corner Mecanum platform \cite{pin1994omnidirectional}, \cite{gfrerrer2008geometry}--\cite{barcenas2025robust} & Planar
holonomic kinematics and omnidirectional motion & Assumes a statically
stable rectangular wheel geometry & Derive the Jacobian from measured
collinear wheel locations and roller signs. \\
Complementary attitude estimation \cite{mahony2008filters} & Combines gyro short-term
response with gravity reference & Acceleration and mounting errors can
bias pitch & Calibrate bias and axes; validate under commanded
acceleration. \\
LQR\slash{}state feedback \cite{astrom2008feedback}, \cite{anderson1990optimal}, \cite{chacko2023trajectory} & Systematic multistate
feedback near an equilibrium & Depends on model, state scaling, and
unsaturated operation & Document weights, discretization, limits, and
robustness tests. \\
CAN-based integrated actuators \cite{espressif2026twai}, \cite{cubemars2026ak4010}, \cite{nxp2026tja1051} & Local FOC
and compact wiring & Network delay, mode settings, encoder location, and
backlash & Measure timing and verify every installed actuator
parameter. \\
\end{longtable}
